% 作者题证的编辑转录副本，不是独立下载的上游原件。
% 来源：https://github.com/OpenLogicProject/OpenLogic/blob/master/content/incompleteness/incompleteness-provability/lob-thm.tex

\paragraph{Notation (editor's expansion of the source macros).}
$\mathrm Q$ denotes Robinson arithmetic. For a formula $A$, $\ulcorner A\urcorner$ is the standard numeral of its G\"odel number; $\operatorname{diag}$ is the primitive recursive diagonal-substitution function. The formula $D_{\operatorname{diag}}(x,y)$ represents this function in $\mathrm Q$.
\begin{lemma}[Fixed-point lemma; the author's support proof]
Let $B(x)$ be any formula with one free variable $x$. Then there is a sentence $A$ such that $\mathrm Q\vdash A\leftrightarrow B(\ulcorner A\urcorner)$.
\end{lemma}
\begin{proof}
Given $B(x)$, let $E(x)$ be the formula $\exists y(D_{\operatorname{diag}}(x,y)\land B(y))$ and let $A$ be its diagonalization, i.e., the formula $E(\ulcorner E(x)\urcorner)$.
Since $D_{\operatorname{diag}}$ represents $\operatorname{diag}$, and $\operatorname{diag}(\ulcorner E(x)\urcorner)=\ulcorner A\urcorner$, $\mathrm Q$ can derive
\begin{align}
&D_{\operatorname{diag}}(\ulcorner E(x)\urcorner,\ulcorner A\urcorner),\label{repdiag1}\\
&\forall y(D_{\operatorname{diag}}(\ulcorner E(x)\urcorner,y)\to y=\ulcorner A\urcorner).\label{repdiag2}
\end{align}
Now we show that $\mathrm Q\vdash A\leftrightarrow B(\ulcorner A\urcorner)$. We argue informally, using just logic and facts derivable in $\mathrm Q$.
First, suppose $A$, i.e., $E(\ulcorner E(x)\urcorner)$. Going back to the definition of $E(x)$, we see that $E(\ulcorner E(x)\urcorner)$ just is
\[\exists y(D_{\operatorname{diag}}(\ulcorner E(x)\urcorner,y)\land B(y)).\]
Consider such a $y$. Since $D_{\operatorname{diag}}(\ulcorner E(x)\urcorner,y)$, by \eqref{repdiag2}, $y=\ulcorner A\urcorner$. So, from $B(y)$ we have $B(\ulcorner A\urcorner)$.
Now suppose $B(\ulcorner A\urcorner)$. By \eqref{repdiag1}, we have
\[D_{\operatorname{diag}}(\ulcorner E(x)\urcorner,\ulcorner A\urcorner)\land B(\ulcorner A\urcorner).\]
It follows that
\[\exists y(D_{\operatorname{diag}}(\ulcorner E(x)\urcorner,y)\land B(y)).\]
But that's just $E(\ulcorner E(x)\urcorner)$, i.e., $A$.
\end{proof}

\paragraph{Derivability conditions (assumptions, from the source).}
For the provability formula $\operatorname{Prov}_T$, assume:
\begin{enumerate}[label=P\arabic*.]
\item If $T\vdash A$, then $T\vdash\operatorname{Prov}_T(\ulcorner A\urcorner)$.
\item $T\vdash\operatorname{Prov}_T(\ulcorner A\to B\urcorner)\to(\operatorname{Prov}_T(\ulcorner A\urcorner)\to\operatorname{Prov}_T(\ulcorner B\urcorner))$.
\item $T\vdash\operatorname{Prov}_T(\ulcorner A\urcorner)\to\operatorname{Prov}_T(\ulcorner\operatorname{Prov}_T(\ulcorner A\urcorner)\urcorner)$.
\end{enumerate}
\begin{theorem}[L\"ob's theorem; the counted problem]
Let $T$ be an axiomatizable theory extending $\mathrm Q$, and suppose $\operatorname{Prov}_T(y)$ is a formula satisfying conditions P1--P3. If $T$ derives $\operatorname{Prov}_T(\ulcorner A\urcorner)\to A$, then in fact $T$ derives $A$.
\end{theorem}
\begin{proof}
Suppose $A$ is a sentence such that $T$ derives $\operatorname{Prov}_T(\ulcorner A\urcorner)\to A$. Let $B(y)$ be the formula $\operatorname{Prov}_T(y)\to A$, and use the fixed-point lemma to find a sentence $D$ such that $T$ derives $D\leftrightarrow B(\ulcorner D\urcorner)$. Then each of the following is derivable in $T$:
\begin{align}
&D\leftrightarrow(\operatorname{Prov}_T(\ulcorner D\urcorner)\to A)\label{L1}\\
&\qquad\text{$D$ is a fixed point of $B(y)$}\notag\\
&D\to(\operatorname{Prov}_T(\ulcorner D\urcorner)\to A)\label{L2}\\
&\qquad\text{from \eqref{L1}}\notag\\
&\operatorname{Prov}_T(\ulcorner D\to(\operatorname{Prov}_T(\ulcorner D\urcorner)\to A)\urcorner)\label{L3}\\
&\qquad\text{from \eqref{L2} by condition P1}\notag\\
&\operatorname{Prov}_T(\ulcorner D\urcorner)\to\operatorname{Prov}_T(\ulcorner\operatorname{Prov}_T(\ulcorner D\urcorner)\to A\urcorner)\label{L4}\\
&\qquad\text{from \eqref{L3} using condition P2}\notag\\
&\operatorname{Prov}_T(\ulcorner D\urcorner)\to\bigl(\operatorname{Prov}_T(\ulcorner\operatorname{Prov}_T(\ulcorner D\urcorner)\urcorner)\to\operatorname{Prov}_T(\ulcorner A\urcorner)\bigr)\label{L5}\\
&\qquad\text{from \eqref{L4} using P2 again}\notag\\
&\operatorname{Prov}_T(\ulcorner D\urcorner)\to\operatorname{Prov}_T(\ulcorner\operatorname{Prov}_T(\ulcorner D\urcorner)\urcorner)\label{L6}\\
&\qquad\text{by derivability condition P3}\notag\\
&\operatorname{Prov}_T(\ulcorner D\urcorner)\to\operatorname{Prov}_T(\ulcorner A\urcorner)\label{L7}\\
&\qquad\text{from \eqref{L5} and \eqref{L6}}\notag\\
&\operatorname{Prov}_T(\ulcorner A\urcorner)\to A\label{L8}\\
&\qquad\text{by assumption of the theorem}\notag\\
&\operatorname{Prov}_T(\ulcorner D\urcorner)\to A\label{L9}\\
&\qquad\text{from \eqref{L7} and \eqref{L8}}\notag\\
&(\operatorname{Prov}_T(\ulcorner D\urcorner)\to A)\to D\label{L10}\\
&\qquad\text{from \eqref{L1}}\notag\\
&D\label{L11}\\
&\qquad\text{from \eqref{L9} and \eqref{L10}}\notag\\
&\operatorname{Prov}_T(\ulcorner D\urcorner)\label{L12}\\
&\qquad\text{from \eqref{L11} by condition P1}\notag\\
&A\qquad\text{from \eqref{L9} and \eqref{L12}.}\notag
\end{align}
\end{proof}
